This chapter surveys the research landscape of tabular data understanding, organized around the two pillars that structure this report: \emph{schema understanding}, which concerns the structural blueprint of tabular data and the mapping of natural-language intent onto it (Section~\ref{sec:schema}); and \emph{data understanding}, which concerns reasoning over the content of tables, canonically instantiated as Table Question Answering (Section~\ref{sec:data}). The two pillars share one underlying difficulty: large language models consume one-dimensional token sequences, whereas a table's meaning is carried jointly by its rows, columns, headers, and the relationships among cells, so information is inevitably lost when a table is linearized for a model~\cite{wu2025tabular}. How a table is \emph{represented} to the model is therefore a decisive variable; because this choice bears most directly on reasoning over content, we treat it within the data-understanding pillar (Section~\ref{sec:representations}) rather than as a separate concern. Rather than cataloguing papers in isolation, we synthesize the literature thematically within each pillar, and we close each pillar by contrasting what current work achieves on academic benchmarks with what it still lacks for real-world, industry-scale deployment, the gap that motivates this report.

%========================================================
\section{Schema Understanding}
\label{sec:schema}
Schema understanding concerns the structural blueprint of tabular data (table and column names, data types, primary and foreign keys, and inter-table relationships) and the task of mapping a user's natural-language intent onto that structure. Its canonical instantiation is Text-to-SQL: translating a natural-language question into an executable query over a relational database. The schema itself is the model's input representation here, expressed for example as SQL \texttt{CREATE TABLE} statements or dataframe declarations that expose structure while omitting the bulk of the values; this sidesteps context-length limits but presupposes that the schema is well-formed and interpretable~\cite{wu2025tabular,sui2024table}. On curated academic schemas Text-to-SQL is largely solved; the difficulty, and the reason the problem remains open, is \emph{industry scale}. Deploying this capability in realistic settings is hard for three reasons: schemas are large (thousands of tables and tens of thousands of columns) and their naming is often cryptic, undocumented, and obsolescent, drifting as the database evolves; the relevant schema elements must be retrieved from among this mass before a query can be formed; and the generated query must be not merely syntactically valid but logically faithful to the user's intent~\cite{li2023bird}. We organize the literature around how it addresses these three challenges.

\subsection{Knowledge-base and metadata construction}
The first challenge (obscure schemas) has motivated a line of work that enriches the raw schema with a semantic layer before any query is attempted. Such knowledge bases are built by aligning schema metadata with business glossaries~\cite{lobo2023matching}, by automatically generating column and table descriptions~\cite{gao2025automatic}, or by injecting domain knowledge from external documents~\cite{ma2024enhancing}. These descriptions disambiguate cryptic identifiers (for example, resolving a column named \texttt{A2} to its business meaning) and externalize schema semantics into a structured artifact that a model can consult, which is essential when industry-scale schemas exceed the model's context window. A recurring limitation of this line of work is that it treats the knowledge base as a \emph{static} asset, constructed once and thereafter fixed; it does not adapt as new queries reveal previously undocumented aspects of the schema.

\subsection{Schema linking and retrieval}
The second challenge (identifying which tables and columns are relevant to a given question) is a retrieval problem. Retrieval-Augmented Generation (RAG)~\cite{lewis2020retrieval} provides the general template of grounding generation in retrieved context, and dense retrievers built on sentence embeddings~\cite{reimers2019sentence} make it possible to match a natural-language query against schema-element descriptions by semantic similarity. For databases, retrieval must respect structure: relevant tables imply related tables through foreign-key relationships, so join-aware and hierarchical retrieval (first selecting tables, then narrowing to columns, and expanding along relationships to preserve join paths) is more effective than flat retrieval. A central difficulty is that vague user queries map poorly onto structured schemas, which has prompted iterative approaches that reformulate the query when the retrieved context is judged insufficient, echoing self-reflective RAG methods developed for open-domain settings~\cite{asai2023self}.

\subsection{From encoder--decoder models to LLM pipelines}
The third challenge (generating faithful queries) has driven the methodological evolution of the field. Early neural systems framed Text-to-SQL as sequence generation: Seq2SQL used reinforcement learning to generate structured queries~\cite{zhong2017seq2sql}, and SQLNet replaced generation with a sketch-based slot-filling formulation to avoid the pitfalls of ordering in the reward~\cite{xu2017sqlnet}. The introduction of large, cross-domain benchmarks such as Spider~\cite{yu2018spider} and, more recently, the industry-scale BIRD benchmark~\cite{li2023bird} shifted evaluation toward generalization across unseen schemas and toward execution accuracy on realistic databases. With LLMs, the dominant paradigm became in-context learning: DIN-SQL decomposes the task into sub-problems and adds self-correction~\cite{pourreza2023din}; DAIL-SQL systematically studies prompt construction and exemplar selection~\cite{gao2023text}; SQLPrompt explores few-shot prompting for the task~\cite{sun2023sqlprompt}; MAC-SQL casts generation as a collaboration among specialized agents (selector, decomposer, refiner)~\cite{wang2024mac}; and CHASE-SQL generates multiple candidate queries via diverse reasoning paths and selects among them with a preference-optimized selector~\cite{pourreza2024chase}. The clear trend is from single-pass generation toward multi-step pipelines that decompose, verify, and correct.

\subsection{Self-reflection for structured generation}
A unifying thread across the most recent methods is \emph{self-reflection}: rather than trusting a single generation, the model scores its own output and iteratively refines it. The idea originates in open-domain generation, where self-reflection improves factuality and consistency by having a model critique and revise its answers~\cite{ji2023towards}, and where self-reflective retrieval decides when and what to retrieve~\cite{asai2023self}. Adapting this paradigm to Text-to-SQL is non-trivial because the domain offers concrete, checkable signals (a query can be executed, its syntax validated, and its results inspected) that open-domain generation lacks. Using an LLM as an evaluator (LLM-as-a-judge) makes it possible to score retrieval relevance, semantic correctness, and result--intent alignment, though such judges must be prompted carefully with rubrics and examples to mitigate the biases inherent in LLM-based evaluation~\cite{zheng2023judging}. Our own work, described in the next chapter, builds directly on this thread by integrating reflection at the retrieval, generation, and verification stages of a Text-to-SQL pipeline.

\subsection{From benchmarks to the enterprise}
Despite rapid progress, the schema-understanding literature is overwhelmingly validated on curated benchmarks such as Spider~\cite{yu2018spider} and BIRD~\cite{li2023bird}, whose schemas are modest in size, consistently named, and fully documented. Three gaps separate these reported results from real-world deployment. First, \emph{scale and obscurity}: production schemas carry thousands of tables and cryptic, undocumented column names, so methods that assume the whole schema fits in context, or that rely on self-explanatory names, degrade sharply. Second, \emph{obsolescence}: real schemas drift as columns are added, repurposed, or abandoned, yet most systems treat the schema, and any semantic layer built over it, as a static artifact fixed at setup time. Third, \emph{the absence of gold queries}: benchmarks supply reference SQL for evaluation, whereas a deployed system must judge the correctness of its own output with no reference to compare against. These gaps motivate the self-reflective, continually-enriched design we present in Chapter~\ref{ch:ch3}.

%========================================================
\section{Data Understanding}
\label{sec:data}
Data understanding concerns reasoning over the \emph{content} of tables (their values and the aggregations, comparisons, trends, and outliers they encode), frequently in combination with surrounding text. Whereas schema understanding privileges the structural blueprint, data understanding privileges the cells themselves and the semantics that connect them. Its canonical instantiation, and the task we concentrate on, is Table Question Answering. We first review how table content is represented to the model, then treat Table QA as the representative task (noting two adjacent tasks for context), and finally consolidate the open challenges and the real-world gap that define our forward direction.

\subsection{Representing table content}
\label{sec:representations}
Before a model can reason over a table, its two-dimensional content must be rendered into a form the model consumes, and performance is sensitive to this choice, varying by several percentage points with the format~\cite{sui2024table}. Three families present content directly, each trading completeness against compactness differently; the schema-based view of Section~\ref{sec:schema} abstracts structure rather than content and belongs with that pillar. \emph{Serialization} renders the table as text (HTML, Markdown, JSON, LaTeX, or delimited variants); it is the default in most benchmarks because such formats occur naturally in pre-training data, but long or multi-table inputs quickly exceed the context window, forcing truncation or row- and column-sampling~\cite{wu2025tabular}. \emph{Image} representations render the table visually, preserving merged cells and hierarchical headers that serialization distorts, at the cost of resolution-bounded legibility as tables grow~\cite{wu2025tabular}. \emph{Table encoders} embed two-dimensional structure directly into the model; table foundation models such as TableGPT2 couple a row- and column-aware encoder with a large base model and outperform generalist prompting across a range of tasks~\cite{su2024tablegpt2}. A limitation common to all three families, and consequential for deployment, is that they are demonstrated on small tables, typically a handful of columns and a few dozen rows, and none scales gracefully to the hundreds of columns and thousands to millions of rows of real-world data.

\subsection{Graph representations}
\label{sec:graph}
A distinct and, we argue, especially promising way to represent a table casts it as a \emph{graph}: cells, rows, headers, and their groupings become typed nodes, and the table's relational topology (which cell lies in which row and column, how headers span sub-columns, and how foreign keys link tables) becomes edges. Unlike the linear formats above, a graph makes a table's structure explicit and \emph{addressable} rather than implicit in a flattened token stream, so a model can navigate directly to the region that matters instead of scanning the entire serialization. This composes naturally with graph-based retrieval, which delivers only the relevant subgraph to the model: GraphRAG builds a graph index for query-focused retrieval~\cite{edge2024graphrag}, G-Retriever retrieves relevant subgraphs for question answering over textual graphs~\cite{he2024gretriever}, and agentic navigators such as GraphOTTER walk a graph of a complex table from headers to the cells that answer a query~\cite{li2025graphotter}. For hierarchical and relationally complex tables, this structure-aware selection is a natural fit for Table Question Answering, where the answer-bearing cells must be located before any reasoning can occur.

The promise, however, is only partly realized. Current graph methods are confined either to \emph{document}- and text-graph settings, where GraphRAG and G-Retriever operate over prose rather than dense tabular data, or to \emph{small, complex} tables, where GraphOTTER demonstrates graph reasoning but scores every traversal step with the LLM, tying latency to table size. Graph representation is therefore a solid foundation whose use today is limited to documentation-style graphs and small-scale tables. The motivation we carry forward is to lift this ceiling: to represent \emph{large} tabular data as a graph and to retrieve over it by addressing pre-computed structure rather than by LLM-scored traversal, so that the structure-awareness that already benefits Table QA on small tables extends to the hundreds of columns and thousands of rows of real-world data. We develop this direction in Chapter~\ref{ch:future}.

\subsection{Table Question Answering}
Table Question Answering (TQA) is the most studied task: given a natural-language question and one or more tables, sometimes accompanied by a passage, the model produces an answer derived from the table content, typically as cell spans, computed values, or short text spans~\cite{wu2025tabular}. The benchmark landscape has expanded along several axes. Foundational datasets such as WikiTableQuestions~\cite{pasupat2015wtq} and the sequential SQA~\cite{iyyer2017sqa} established open-domain answering over web tables. Subsequent benchmarks pushed on \emph{domain} (finance in TAT-QA~\cite{zhu2021tatqa} and FinQA~\cite{chen2021finqa}), on \emph{input complexity} (hybrid tabular--textual reasoning in HybridQA~\cite{chen2020hybridqa}, hierarchical tables in HiTab~\cite{cheng2022hitab}, and multi-hierarchical and multi-table settings in MULTIHIERTT~\cite{zhao2022multihiertt} and MultiTabQA~\cite{pal2023multitabqa}), and on \emph{answer format} (free-form generation in FeTaQA~\cite{nan2021fetaqa}). A notable methodological finding is that decomposition-based pipelines, which break a question into simpler steps executed against the table, have driven large gains and now approach saturation on the simpler retrieval-oriented benchmarks. A caveat qualifies these gains: the tables in these benchmarks are small, and the decomposition pipelines that drive the gains issue an LLM call per step, so both the evaluation and the method strain once tables reach industry scale, where the answer-bearing cells are a vanishing fraction of the content. This foreshadows the robustness and efficiency challenges discussed below.

\subsection{Adjacent tasks}
Two further tasks round out the data-understanding pillar but are not our focus. \emph{Table Fact Verification} (TFV), also called table-based natural language inference, asks whether a claim is supported, refuted, or unverifiable given a reference table; TabFact established the task at scale over Wikipedia tables~\cite{chen2020tabfact}, and modern methods exceed 80\% accuracy, though the strong numbers concentrate on well-structured single tables and thus test fact retrieval more than deep reasoning. \emph{Table-to-Text} generation instead asks a model to describe table content faithfully in natural language, ranging from controlled descriptions of highlighted cells (ToTTo~\cite{parikh2020totto}) to query-focused summaries that must reason about underlying causes rather than surface content (QTSUMM~\cite{zhao2023qtsumm}). Both tasks share Table Question Answering's core difficulty, isolating and reasoning over the relevant cells, which is why we take Table QA as the representative task for the pillar.

\subsection{Open challenges}
Synthesizing across tasks, the survey literature identifies three challenges that define the frontier of data understanding~\cite{wu2025tabular}.

\emph{First, limited scope beyond retrieval and arithmetic.} Many benchmarks are constructed by translating SQL queries or mathematical expressions into natural-language questions, so their core structure rewards retrieval and computation rather than higher-order reasoning. As several such benchmarks approach saturation, recent work has begun to demand diagnostic, predictive, and prescriptive thinking (identifying insights, forecasting, and generating charts from ambiguous requests) and intent-driven querying, where a system must infer what an abstract, goal-oriented request actually asks for~\cite{lei2024spider2,he2024text2analysis}. These tasks cannot be solved by matching synonyms to table headers; they require understanding the table's overall context and the user's intent.

\emph{Second, lack of robustness on complex inputs.} Intricate table structures, long tables, lengthy accompanying texts, and multi-table settings degrade model performance far more than they degrade human performance. On hierarchical and multi-table benchmarks such as HiTab~\cite{cheng2022hitab} and MULTIHIERTT~\cite{zhao2022multihiertt}, reported accuracies fall below 50\% against human performance around 83\%, and even on well-structured multi-table QA the strongest evaluated models reach only about half of human accuracy. Robustness to input complexity is thus a first-order open problem.

\emph{Third, limited generalization across representations.} As noted in Section~\ref{sec:representations}, performance varies by up to several percentage points with the input format, and models struggle to transfer across the diverse representations used by different benchmarks~\cite{sui2024table}. Proposed remedies include standardizing serialization options, deriving serialized views from images of complex tables, and training on serialization-to-serialization tasks that convert one format into another to build robustness to varied input styles.

\subsection{From benchmarks to the enterprise}
These challenges share a root cause that parallels the schema pillar: data-understanding methods are developed and scored on small, clean, well-structured tables, whereas real applications demand reasoning over data that is large, noisy, hierarchical, and spread across many tables. The field is near saturation on simple retrieval-oriented benchmarks yet remains far below human performance the moment inputs resemble production data, and the pipelines that post the strongest numbers scale their cost with table size, so they do not transfer to tables that cannot fit in a model's context. Reasoning faithfully over content at real-world scale and complexity, rather than posting incremental gains on already-saturated simple benchmarks, is the frontier this report takes as its forward direction.

Having mapped both pillars and the real-world gaps that remain, the next chapter turns to our own contribution to the schema-understanding pillar, Reflect-SQL, which operationalizes this analysis and lays the groundwork for pursuing data understanding at industry scale.
